\section*{Methods}
\label{sec:method}

%%% NOTES. (journal instructions)
%%%
%%% The Methods should include detailed text describing any steps or procedures used in producing the data, including full descriptions of the experimental design, data acquisition assays, and any computational processing (e.g. normalization, image feature extraction). See the detailed section in our submission guidelines for advice on writing a transparent and reproducible methods section. Related methods should be grouped under corresponding subheadings where possible, and methods should be described in enough detail to allow other researchers to interpret and repeat, if required, the full study. Specific data outputs should be explicitly referenced via data citation (see Data Records and Citing Data, below). Authors should cite previous descriptions of the methods under use, but ideally the method descriptions should be complete enough for others to understand and reproduce the methods and processing steps without referring to associated publications. There is no limit to the length of the Methods section. Subheadings should not be numbered.


To create a dataset of text reuse in scientific texts, four things are needed: an operationalization of the phenomenon of text reuse, a large collection of scientific publications, detailed metadata for each of these publications, and a scalable yet effective method for detecting reuse in these publications. In this section we describe all four, explaining the latter both generally and formally to ensure their reproducibility.
%First, we define our operationalization of text reuse in the context of this paper. Then, a general overview on the process of detecting reused text in large document collections is given, followed by detailed descriptions of our data acquisition and preprocessing, as well as the individual steps taken in calculating all document-to-document comparisons to identify reused text between them.
Figure~\ref{fig:pipeline} illustrates the processing pipeline for the creation of the dataset. 

\subsection*{An Operationalization of Text Reuse}

To reuse something means to use it again after the first time. Reused text is text that is primarily, if not exclusively, derived from another text. In academic writing, writing techniques for reusing a text include boilerplate, quotation, paraphrasing,
%
%\footnote{Paraphrasing is also central to translation \cite{baker:2020}: it goes back to the translation practice of Greek and Roman texts, in which two types of translations can be distinguished: Word-for-word translations (or metaphrases) and sense-for-sense translations (or paraphrases). These translation techniques are ``faithful''  to the original, as opposed to ``unfaithful'' free translations, although this dichotomy is controversial. We consider only intra-linguistic reuse, where paraphrasing is common and metaphrasing is equivalent to quotation.}
and summarizing. What all these techniques have in common is that a reused text and its original have a certain kind of similarity that the reader can recognize \cite{weber-wulff:2002}. Manual identification of text reuse between two given texts is therefore based on the identification of the relevant text passages where such similarities can be detected.

%%% ADDTION. (more precise, but also something we did not do in the end; keep for later)
%Manual detection of text reuse is therefore based on identifying parts of two texts that have unlikely similarities. The less likely it is that a particular form of similarity occurs in two independently written texts, and the more such similarities are identifiable, the more certain it is to assume that one of the texts was actually reused to write the other.

To operationalize this notion of reuse, a suitable text similarity measure and a similarity threshold have to be chosen, where a pair of texts has to exceed the threshold in order to be considered potential reuse. We broadly distinguish similarity measures that operate at the syntactic level from measures that operate at the semantic level of language, where the former capture the ``reuse of words'' and the latter the ``reuse of ideas'' \cite{sadeghi:2019}.
%%% ADDITION.
%\bscom{Considering reuse, this distinction reflects the age-old distinction between ``reuse of words'' and ``reuse of ideas'', outlined as early as the thirteenth century~AD by Persian scholar Shams-e-Qays \protect\cite{sadeghi:2019}. Widely adopted today \protect\cite{martin:2015,barron:2013}, sometimes both kinds appear in conjunction \protect\cite{moskovitz:2021}.}{All of this leads too far away...}
In practice, text reuse detection relies heavily on syntactic similarity detection. \cite{moskovitz:2021,vroniplag:2021}. Syntactic reuse is relatively easy to visualize and consequently faster to check than semantic reuse \cite{riehmann:2015}. The former can be checked by distant reading, the latter requires close reading, which causes high costs with increasing text length \cite{moretti:2013}. Similarly, it can be assumed that semantic reuse is much less common than syntactic reuse, since the most common goal of text reuse is to save time and cognitive effort, whereas the time savings of semantic reuse are generally lower \cite{potthast:2013c}.

We opt for a consvervative similarity analysis at the syntactic level, measuring the correspondence between the surface forms of words occurring in two given texts and the phrases formed from them. This design decision is also motivated by the target domain: In addition to the contribution of new ideas, a large fraction of scientific contributions describe reflection on as well as advances to known ideas, and the development of solutions to tasks and problems up to the point of transfer to practice. A semantic similarity score in this context would rather lead to a citation graph mixing natural matches of ideas with intended reuse. It is in the nature of current semantic similarity measures that they often capture even minor and distant similarities, leading to significant noise in the form of false positives in tasks such as text reuse detection. Also, detection at the syntactic level is more in line with the concept of text reuse used in practice.In other words, `similar' here refers to an overlap in the vocabulary and wording of two documents, an operationalization common to large-scale text reuse detection \cite{martin:2015}.

Although even a single word overlap between two pieces of text, such as a specific and unlikely spelling mistake, can be sufficient as a strong indication of reuse, automatic methods are not yet able to reliably detect such cases. Instead, several overlaps of words and phrases are required that occur in close proximity to each other in both texts, thus forming potentially reused text passages. Such passages are not necessarily verbatim copies: Words and phrases may be added, removed, or changed, and sentences may be rearranged. Nevertheless, sufficient overlap must remain, which can be modeled by the similarity threshold parameter mentioned above. This parameter cannot be derived formally, but must be determined empirically.

This operationalization of text reuse, and thus the cases included in our dataset, are orthogonal to the question of the legitimacy of reuse. Thus, the notion of ``text reuse'' is explicitly not limited to or equated with ``plagiarism.'' We do not draw any conclusions from the fact that two documents have a reused passage in common, but merely observe that there is one. The distinction between original and plagiarized text cannot and should not be an algorithmic decision \cite{weber-wulff:2019}. Further limitations are discussed in the Usage Notes section.


%Limitations:

% - A certain kind of reuse falls into the broader realm of intertextuality, such as when an author draws inspiration from or alludes to another's writing through analogy, perhaps even chosen to be understood only by an insider group. Although contextual information can still serve as evidence for this type of reuse, we omit such artful cases of reuse from our operationalization. Unlike in literary writing, for example, we assume that they are rare in the vast majority of scientific writing as a whole.

% - We also do not identify which of the two texts is the source and which the reused version, nor whether an unseen third text is the actual source, as this contextual data, despite having at hand the metadata of the paper, is insuficient to deduce this property.


\subsection*{Text Reuse Detection in Large Document Collections}

Detecting all reuses of text in a collection of documents is a problem of high computational complexity. Each document must be compared pairwise with every other document in the collection. So for $N$~documents, $N\cdot (N-1)$~document-to-document comparisons must be made. Each individual comparison requires a considerable amount of time and resources, so that the processing of larger collections of documents exceeds even large computational capacities. In order to still be able to analyze large document collections for text reuse, the set of comparisons is pruned by filtering out document pairs that are guaranteed not to exhibit text reuse according to our operationalization.

This is achieved through a two-step process: first, a cheap-to-compute heuristic is applied to identify candidate pairs of documents where text reuse is likely. Then, the expensive document comparison is performed only for the candidate pairs identified in this way. All other pairs are skipped.The first step of this two-part process is commonly called a source retrieval \cite{stein:2007d}, while the second step is called text alignment \cite{potthast:2013e}. The result of the whole process is a set of cases of text reuse between documents in the collection. A case of text reuse is modeled here as a pair of text passages, one in each document involved in the comparison, that have sufficient overlap of words and phrases, along with references to their source documents and where exactly they are found in them.

While source retrieval greatly improves the efficiency of the overall process, it introduces an error in the form of reduced recall, i.e., pairs of reused passages are overlooked that would have been discovered by an exhaustive comparison of each pair of documents. The goal of source retrieval is therefore to filter out only those document pairs for which the text alignment step will not find any reused passages, which is to be expected for the vast majority of all document pairs. We choose the source retrieval parameters conservatively, so that only those document pairs are skipped for which the text alignment is guaranteed to return no result.

%%% ADDITION.
% Altogether, to detect reused text passages in a document collection~$D$, we retrieve for each document $d \in D$ the set $D_d$ of candidate documents and compute for each $d'\in D_d$ the set of reused text passages $R_{d,d'}$. The union of all detected passages is used as an approximation of~$R_D$. The goal of the source retrieval step is to approximate the recall of an exhaustive comparison while minimizing the sum of the sizes of retrieved candidate documents $\sum_{d\in D} |D_d|$, i.e., the computing capacity actually expended. 

\subsection*{Publication Acquisition, Preprocessing and Metadata}\label{sub:data}

Detecting scientific text reuse and contextualizing it on a large scale requires a large collection of scientific publications and detailed metadata about them. We compiled such a dataset in the five steps of document selection, plain text extraction, text preprocessing, metadata acquisition, and metadata standardization.

We build on the CORE dataset \cite{knoth:2012}, one of the largest collections of open-access scientific publications sourced from more than 12,000~data providers. First, we identify the 6,015,512~unique open-access~DOIs in the March~1,~2018 CORE dataset. Since the plain texts extracted from the PDF files of the publications, as provided by the CORE data, are of varying quality and no structural annotations (such as markup for citations, in-text references, section annotations) are available, we chose to obtain the original PDF~files of the identified open-access~DOIs from various publicly available repositories.

The plain text extraction has been repeated on the acquired PDF~files using the standardized state-of-the-art toolchain GROBID \cite{lopez:2015}. A minimum of~1,000 and a maximum of 60,000~space-separated words are introduced as an effective heuristic to filter out common plaintext extraction errors. In total, we obtained and extracted clean plaintext for 4,267,166~documents (70\%~of the open access publications in the original CORE~dataset). Since text alignment analyzes word overlaps between documents, the highly standardized meta-information in scientific texts such as citations, numbers, author names, affiliations, and references lead to exponentially more false-positive detections of reuse than without them. Therefore, text alignment processes the abstract and main body of a document without in-text references, tables, figures, bibliographic data, and numeric data, while normalizing or removing special characters, which minimizes the false positives observed in preliminary studies. However, research on text reuse and its evaluation on a case-by-case basis benefits from, or even requires, the in-text metadata mentioned above, so that two versions of each text passage involved in a reuse case were kept: the one that came from GROBID, including the aforementioned information except for tables and images, and, the one that was fed into the text alignment.

Based on the detections, we augment the metadata provided by CORE with additional data from the Microsoft Open Academic Graph~(OAG) \cite{tang:2008,sinha:2015}, which contains study field annotations for a large number of publications. Metadata is assigned by matching records in~CORE and~OAG using an article's DOI~identifier. Since the annotated disciplines in the~OAG do not follow a hierarchical scheme and since they are of different granularity per publication (e.g., ``humanities'' as a whole vs. ``chemical solid state research'' as a subfield of chemistry), we manually map the classification found in the~OAG to the standard hierarchical \textit{DFG~Classification of Scientific Disciplines, Research Areas, Review Boards and Subject Areas} \cite{dfg:2016}.%
\footnote{We have chosen to replace the term ``review board'' used in the DFG~classification with the more conventional term ``field of study''.}
The mapping was done by three people independently. In the few cases where there was disagreement, consensus was reached through discussion. The final mapping includes 46~individual scientific fields of study drawn from 14~scientific areas and four major scientific disciplines.

\subsection*{Source Retrieval}

An exhaustive comparison of all four million publications greatly exceeds the available computing capacity. Therefore, a source retrieval step has to be carried out first \cite{hagen:2015e,hagen:2017e}. By applying suitable heuristics, the source retrieval step can be operationalized with linear time complexity with respect to the document count \cite{stein:2019c}, compared to the quadratic expenditure of an exhaustive comparison.

The source retrieval component of our pipeline, i.e., the computation of all candidate comparisons~$D_d$ for a given document~$d$, is operationalized by treating text reuse as a locally bounded phenomenon. Candidate pairs are presumed to be identifiable by comparing text passages between two documents. If the similarity is sufficiently high for at least one combination of passages between two documents~$d$ and~$d'$, then~$d'$ is considered a candidate for~$d$. 

To achieve high scalability, we implement the similarity computation by applying locality-sensitive hash functions~$h$ to each passage~$t$ in a document~$d$, thus representing each document's passages with a set of hash values~$h(t)$. These can be interpreted as a set of fingerprints, encoding the characteristics of each passage in a document in compressed form. The similarity between two passages~$t$ and~$t'$ stemming from two different documents~$d$ and~$d'$ is then approximated by the extent of overlap between their hash sets,~$|h(t)\cap h(t')|$. Thus, a document~$d'$ is considered a candidate source of text reuse for a document~$d$, if at least one of their passage-level hash sets intersects \cite{stein:2019c}: $\exists t \subseteq d$ and $\exists t' \subseteq d' : h(t) \cap h(t')\neq\emptyset$. This means that two documents have to share at least one of the passage-level fingerprint hashes, i.e. have at least one common characteristic, to be deemed a candidate pair.

In our pipeline, the documents are first divided into consecutive passages of $n$~terms, which in turn are each embedded using a bag-of-words representation, i.e., as their term occurrence vector. To find matching passages between two documents, MinHash \cite{broder:97} is chosen as hashing scheme for~$h$. For each passage vector, multiple individual hash functions are applied to each element, and the minimum hash value of each pass is saved, yielding the set minimum hashes for a passage. It can be shown that, for~$m$ individual hashes per passage, two texts with a Jaccard similarity of at least~$m^{-1}$ are guaranteed to produce a hash collision between their hash sets, rendering them a suitable approximation for the lower bound of document similarity \cite{broder:97}. Applying this scheme to all passages in all documents, we detect all document pairs that are locally similar, i.e., which have a word overlap in their bag-of-words passage representations. For example, if documents would be split into~$n=20$ word passages, and~$m=5$ hash fingerprints are calculated per passage, two documents would be identified as pair if they share at least $4$~words across one of the respective passages. A document is never compared to itself.

Hash-based source retrieval allows for a significant reduction of the required computation time, since the hash-based approximation has a linear time complexity with respect to~$|D|$, as opposed to the quadratic complexity of vector comparisons \cite{stein:2019c}. As a result, our source retrieval computation time could be fitted into the allotted budget of two months of computing time on a 130-node Apache Spark cluster, with 12~CPU cores and 196~GB~RAM per node (1560~cores, 250~TB~RAM total). Besides its efficiency, the outlined approach provides us with a second highly beneficial characteristic---the source retrieval step depends on only two parameters: the passage length~$n$, which can be chosen according to computational constraints, and the number of hashes~$m$, which is chosen according to the required minimum similarity separating two passages. By choosing this bound of $m^{-1}$~Jaccard similarity lower than the minimum detection threshold of the subsequent text alignment step, the search space is pruned without a loss in accuracy. Thus, our hash-based similarity search primarily impacts the precision, but not the recall of the source retrieval step with respect to the subsequent text alignment step.


\subsection*{Text Alignment}\label{sub:alignment}

Based on the reduced set of comparison candidates~$D_d$ obtained by source retrieval, the text alignment component extracts the exact location of the reused passages of each candidate pair of documents~$d$ and~$d'\ in D_d$. Here we follow the seed-and-extend approach to local sequence alignment \cite{potthast:2012b}. The two-step process is illustrated in Figure~\ref{fig:alignment}. First, both texts are divided into small chunks (\emph{`chunking'}). Then, the matching chunks in the Cartesian product of the two sets of chunks are computed according to a similarity function~$\varphi$. Its purpose is to identify matching textual chunks (`\textit{seeds}`') between two documents~$d$ and~$d'$ that have the same meaning or can be considered instances of the same concept. Finally, sufficiently close matches are combined into larger passages (\emph{`extension'}). Such a pair of passages~$(t \subseteq d, t' \subseteq d')$ in both documents is then output probable reuse case.

With respect to the chunking and seeding steps, two methodological decisions must be made. First, how to divide a given document into small chunks. And second, which similarity function should be used to compare pairs of text chunks to determine whether they have the same meaning. A widely used approach to the former, which has been shown to produce accurate results, is to divide a text into (word)~\emph{n}-grams, i.e., contiguous chunks of words of length~\emph{n} \cite{stamatatos:2011,citron:2015}. These \emph{n}-grams are overlapping and are created by ``sliding'' a window over the text. For example, the sentence ``The quick brown fox jumps over the lazy dog.'' consists of the \emph{4}-grams ``The quick brown fox'', ``quick brown fox jumps'', ``brown fox jumps over'', ``fox jumps over the'', ``jumps over the lazy'', and ``over the lazy dog'' at an overlap of three words. The  \emph{n}-gram chunks of the two documents~$d$ and~$d'$ are then exhaustively compared, i.e., every part of~$d$ is compared to every part of~$d'$. As a similarity function, we again use a hash function, but a string hash function rather than a locale-sensitive function. This choice enables a linear-time comparison of the Cartesian product. Two parameters can be modified in this approach: the \emph{n}-gram size~$n$, and the \emph{n}-gram overlap~$k$.

The matching chunks found through hash collisions indicate ``seeds'' of a potentially longer case of reused text. To determine whether such a case can be found, the \emph{extension} step joins co-aligned matching chunks into longer passages when a sufficient number are found close to each other. In this manner, not only cases of coherent reuse, but also cases where text was copied and then paraphrased can be detected. For example, if consecutive sentences in a source text are copied into a target text, and then another (new) sentence is placed in between them, this would be still considered a case of coherent reuse. Seeding would succeed in identifying the two copied sentences on their own, yet the extension recognizes that both seeds are in close proximity in both documents, and thus outputs a single case of reuse including both. The opposite case is also possible: two chunks from different locations in a source text are placed close to each other in the target texts. Here, again, an extension is required to reconstruct the full scope of reuse. Note that the extension approach depends on a single parameter,~$\Delta$, which is the maximum distance in characters for two seeds in either of a pair of documents below which a single reuse case is assumed \cite{stamatatos:2011}.

We supply a massively parallelizeable implementation of both the source retrieval and text alignment step, which allows for detecting text reuse in the highly scalable manner needed given the amount and length of input documents: 4.2~million publications, totaling 1.1~terabytes of text data. Overall, the text alignment step accounted for an additional 3.5~months of computing time on the aforementioned Spark cluster.

